\section{Joint horizon and accuracy in description length}\label{sec:jointdescription68}
The exponential test in Lemma~\ref{lem:choitesting67} resolves pairs at
fixed error, but its lower bound is vacuous when the two experiments are
very close. A bounded statistic gives the local estimate needed to retain
the accuracy parameter. Throughout this section $N\ge1$ and the adaptive
distance $d_N$ is the unhalved final-state trace distance maximized over
one common tester, as in \eqref{eq:adaptivemetric67}. The unknown device
is the same memoryless instrument at every call; commands are classical,
and stopping is public and bounded by $N$.

\begin{lemma}[A local binary test]\label{lem:binarylocal68}
For $0\le p,q\le1$,
\begin{equation}\label{eq:binarylocal68}
 \TV\bigl(\operatorname{Bern}(p)^{\otimes N},
          \operatorname{Bern}(q)^{\otimes N}\bigr)
 \ge \frac3{32}\min\{\sqrt N\,|p-q|,1\}.
\end{equation}
The estimate is uniform up to the endpoints of the probability interval.
\end{lemma}
\begin{proof}
Assume $p<q$ and first suppose $q-p\le N^{-1/2}$. Put
$b=N(p+q)/2$ and let $f:\R\to[0,1]$ be the nondecreasing function
which is affine with slope $(8\sqrt N)^{-1}$ on
$[b-4\sqrt N,b+4\sqrt N]$, zero to its left, and one to its right.
For $S\sim\operatorname{Bin}(N,t)$, differentiating the finite binomial
polynomial gives the exact identity
\[
 \frac d{dt}\E_t f(S)
 =N\E_t\{f(T+1)-f(T)\},\qquad T\sim\operatorname{Bin}(N-1,t).
\]
It holds on $(0,1)$ and extends continuously to both endpoints. For
$t\in[p,q]$, the distance from $\E_tT$ to $b$ is at most
$\sqrt N/2+1$. Also $\operatorname{Var}_tT\le(N-1)/4$. Chebyshev's
inequality gives
\[
 \Pr_t\{|T-\E_tT|\le\sqrt N\}\ge3/4.
\]
On this event, both $T$ and $T+1$ belong to the affine interval:
their distance to $b$ is at most $3\sqrt N/2+2\le4\sqrt N$.
Since $f$ is nondecreasing everywhere,
\[
 \frac d{dt}\E_t f(S)\ge\frac3{32}\sqrt N.
\]
Integration from $p$ to $q$, and the characterization of total
variation by $[0,1]$-valued statistics, prove the assertion in this
case. The statistic depends only on the count, which is a function
of the full binary string.

If $q-p>N^{-1/2}$, apply the same construction to
$p$ and $q_0=p+N^{-1/2}<q$. Couple Bernoulli samples by common
uniform random variables. Monotonicity of $f$ then gives
$\E_qf(S)\ge\E_{q_0}f(S)$, and the preceding lower bound is $3/32$.
The cases $p=q$ and $q<p$ follow by equality and symmetry.
\end{proof}
The proof uses only binomial differentiation and a bounded statistic;
it requires neither a normal approximation nor a lower variance bound.
In particular its constant does not deteriorate for a rare outcome.

\begin{corollary}[Local repeated Choi separation]\label{cor:localchoi68}
For the instrument body $\mathcal C_{d,n,m}$, let $t=|J-L|_2$.
Then
\begin{equation}\label{eq:localchoi68}
 d_N(J,L)\ge\frac3{16}
                  \min\left\{\frac{\sqrt N}{2d}t,1\right\}.
\end{equation}
\end{corollary}
\begin{proof}
Use the same normalized Choi inputs as in
Lemma~\ref{lem:choitesting67}. The positive spectral projection of
$d^{-1}\bigoplus_y(J_y-L_y)$ defines a binary test with probability
gap at least $t/(2d)$. Repeating this test independently and applying
Lemma~\ref{lem:binarylocal68} gives \eqref{eq:localchoi68}; classical
trace distance is twice total variation. This is one admissible
nonadaptive tester. For each pair it is the same tester in both
experiments, and it is allowed in the supremum defining $d_N$.
\end{proof}

\begin{theorem}[The joint interior coding law]\label{thm:jointcode68}
Fix $d,n,m$ with $s=d^2(mn^2-1)>0$ and $0<a<1/(mn)$.
For all $N\ge1$ and $0<\delta\le1/32$,
\begin{equation}\label{eq:jointcode68}
 \max\left\{1,\left(\frac{R\sqrt N}{32d\delta}\right)^s\right\}
 \le\mathcal M_N(\delta;a)
 \le\mathcal M_N^{\Q}(\delta;a)\le(2B_N+1)^s,
\end{equation}
where $R=1/(mn)-a$ and $B_N$ is the explicit grid in
Theorem~\ref{thm:adaptiveentropy67}. Consequently the optimal legal
fixed-length description, with or without rational decoded matrices,
has length
\begin{equation}\label{eq:jointbits68}
 \frac s2\log_2N+s\log_2(1/\delta)+O_{d,n,m,a}(1),
\end{equation}
uniformly on this whole range of $N$ and $\delta$. The decoded
instrument is reused unchanged, and no spectral-gap hypothesis occurs.
\end{theorem}
\begin{proof}
The tuple-Frobenius ball of radius $R$ about $J^\circ$ in $V$ is
contained in $\mathcal C_a$. Take a maximal set in this ball whose
pairwise Euclidean distances are greater than
$h=32d\delta/\sqrt N$. Maximality and $s$-dimensional volume give at
least $(R/h)^s$ points. Corollary~\ref{cor:localchoi68} separates
each pair by more than $2\delta$ in $d_N$: before the minimum saturates
its value is greater than $3\delta$, and after saturation it is
$3/16>2\delta$. Thus no legal centre of a $d_N$-ball of radius
$\delta$ can cover two packing points, even if that centre has no
positive eigenvalue margin. This proves the lower bound.

The intrinsic rational codec and its classical-programme estimate in
Theorem~\ref{thm:adaptiveentropy67} give the upper bound with the
same $B_N$. Since $N\ge1$ and $\delta\le1/32$, there is a constant
$C=C(d,n,m,a)$ such that $2B_N+1\le C\sqrt N/\delta$. The lower
bound and this inequality give \eqref{eq:jointbits68}, including
rounding a logarithm up to a fixed-length binary code. When the
volumetric lower expression is less than one, its logarithm is negative;
the nonnegative code length still gives the same uniform additive
estimate. Neither bound charges the queries of a learning procedure.
\end{proof}
The earlier fixed-error theorem is retained. The new assertion is the
uniform additive remainder in \eqref{eq:jointbits68}, which is
independent of both the horizon and the accuracy. It does not assert
sharp dimension-dependent constants, nor is it a lower bound for the
mutable workspace of the variable-description trajectory algorithm.
